% =====================================================================
%  Harald Høffding, FILOSOFISKE PROBLEMER.  Indbydelsesskrift til
%  Kjøbenhavns Universitets Aarsfest til Erindring om Kirkens Reformation,
%  November 1902.  Kjøbenhavn: Universitetsbogtrykkeriet (J. H. Schultz),
%  1902.  90 pp.
%  TRANSCRIPTION (Danish) of the 1902 print, complete.
%
%  SOURCE: the Royal Library's (KB) scan, ~/bibliotek/Høffding, Harald/
%  1902-filosofiske-problemer.pdf (SCANS.tsv).  PAGE MAP (pagemap.py):
%  printed page = PDF page - 13; p. 1 = PDF 14 ... p. 90 = PDF 103.  Front
%  matter (unnumbered): the Indbydelsesskrift leaf PDF 4 (the same leaf
%  again at PDF 6), title PDF 8, preface PDF 10, contents PDF 12; these
%  carry "% --- front matter ---" comments and no \opage.  Not
%  transcribed: KB's cover sheet (PDF 1), blank and library leaves, the
%  bound-in university list after p. 90 (PDF 104-123).  Every page turn
%  is marked "% --- p. N ---" + \opage{N} at the first word of the printed
%  page; a word divided over a page turn ends in %.
%
%  CONVENTIONS: italic -> \textit; small capitals -> \textsc (normal case
%  inside); printed dashes -> ---, number ranges -> --; quotation marks
%  „…“ as printed; subscripts as printed (pp. 50-51).  Høffding's notes
%  are endnotes („Noter“, pp. 85-90, in small type), each headed by its
%  label and the page it belongs to; in the text the anchors are a
%  superscript figure and a parenthesis („fremgaaet“²⁸)“).  The labels run
%  1-12, 12b, 13-31, 31b, 31c, 32-60, as printed; in the Noter the 12b note
%  is labelled „12)“ (marked PRINTED AS IS).  Not transcribed: folios,
%  sheet signatures („12“, p. 89), library stamps.
%
%  METHOD (the e-book method, 2026-10-02; TRANSCRIPTION-PLAYBOOK §00 point
%  0).  No one typed the text.  The base text -- words, punctuation,
%  paragraphs, italics, small caps, headings, note texts -- is the SAGA
%  e-book (Filosofiske_problemer.epub, in this folder); it was aligned
%  letter by letter against the scan's own text layer by
%  ebook-method/filosofiske-problemer/build.py (tracked, repo root), which
%  also takes every page turn from the layer.
%  Of the letter and figure disagreements, the script settled by rule
%  those that are the layer's known faults (glyph confusions, letterspaced
%  small capitals, dropped accents, specks, the note labels); every other
%  one was decided at the image from crop sheets (sheets.py, same folder):
%  52 in the text, 62 in the notes, 22 punctuation disagreements in the
%  text and 30 in the notes (among them every comma closing an italic
%  title and every note's last mark), and 4 paragraph disagreements
%  (decisions.tsv, patches.tsv).  Spot pages 9 and 87 read in full
%  against the image afterwards; p. 87 led to the rule that a layer glyph
%  fault settles a disagreement only where the e-book's word is attested
%  elsewhere in the repository.  The print corrects the
%  e-book at: p. 2 „dér“, „énhver“ (accents); p. 14 Aarsagsforholdet;
%  p. 17 Flechsigs; p. 22 Sammenhøren; p. 23 saadant; p. 41 Genstandes;
%  p. 42 Molekuler; p. 46 a stray point removed; pp. 50-51 the
%  subscripts; p. 54 fører; p. 64 hele, Verdensprincipet; p. 65 føres;
%  p. 68 Resultat; p. 70 har; p. 71 Existens; p. 72 Ret, Principet;
%  p. 77 Friedrich, and „som Selvhævdelsen“, which the e-book omits; in
%  the Noter: Der, Roberto, apparet, p. 36, er, Bevidsthedsfænomenerne,
%  Psykologiske, experience, nur, nicht, Wikner, her, Grundlag,
%  sufficiunt, Wissenschaftslehre, „p. 32 f.“, the point (not comma)
%  after nine italic titles, the Greek of note 60, the note labels.
%  PRINTED AS IS: p. 47 „1dentitetsforhold“ (the figure 1 for I);
%  „(S 3)“, „(S 8)“, „(S 46)“ in the Noter; O with
%  subscript s for O-zero, pp. 50 and 51; the second „12)“.  One doubtful
%  point (p. 89, after „Hegel“) is marked % UNSURE.
%
%  RESIDUAL RISKS: (1) a misprint that the e-book and the layer both
%  normalise is invisible to the method; (2) italics and small capitals
%  are the e-book's, which the layer cannot witness; (3) punctuation is
%  witnessed only where the layer reads it, and quotation marks and
%  dashes come from the e-book alone; (4) letterspacing, if any, is not
%  recorded.
% =====================================================================
\documentclass[12pt,a4paper,openany]{book}
\usepackage[utf8]{inputenc}
\usepackage[T1]{fontenc}
\usepackage{textalpha}
\usepackage{libertinus}
\usepackage{libertinust1math}
\usepackage{graphicx}    % for \reflectbox (the old Danish ɔ: id-est mark)
\DeclareUnicodeCharacter{0254}{\reflectbox{c}}  % ɔ (U+0254) = reversed c
\usepackage[danish]{babel}
\usepackage[protrusion=true,expansion=true]{microtype}
\usepackage{setspace}
\usepackage[margin=1.2in]{geometry}
\usepackage{fancyhdr}
\setlength{\headheight}{28pt}
\addtolength{\topmargin}{-13.5pt}
\usepackage{hyperref}
\hypersetup{pdftitle={Filosofiske Problemer}, pdfauthor={Harald Høffding}, hidelinks}
\pagestyle{fancy}
\fancyhf{}
\fancyhead[LE,RO]{\thepage}
\fancyhead[RE]{\textit{Harald Høffding: Filosofiske Problemer}}
\fancyhead[LO]{\textit{1902}}
\setstretch{1.3}
\title{\textbf{Filosofiske Problemer}}
\author{Harald Høffding}
\date{Kjøbenhavn: Universitetsbogtrykkeriet (J. H. Schultz), 1902}
\usepackage{marginnote}
\newcommand{\opage}[1]{\marginnote{\footnotesize\textit{[#1]}}}
% Contents lines (front matter).
\newcommand{\tocline}[3]{\noindent\makebox[2.2em][l]{#1}#2\dotfill\ #3\par}
\newcommand{\tocsub}[2]{\noindent\hangindent=4.4em\hspace*{2.2em}#1\dotfill\ #2\par}
\newcommand{\tocplain}[1]{\noindent\hspace*{2.2em}#1\par}

\begin{document}
\maketitle
\thispagestyle{empty}
\clearpage

% --- front matter: Indbydelsesskrift (PDF 4; the same leaf again at PDF 6) ---
\begin{center}
INDBYDELSESSKRIFT\\[2pt]
{\scriptsize TIL}\\[2pt]
{\large KJØBENHAVNS UNIVERSITETS}\\
{\large AARSFEST}\\[2pt]
{\scriptsize TIL ERINDRING OM}\\[2pt]
{\large KIRKENS REFORMATION}\\[2pt]
NOVEMBER 1902\\[2em]
% (the University's seal)
HARALD HØFFDING:\\
FILOSOFISKE PROBLEMER\\[1em]
\rule{1cm}{0.4pt}\\[1em]
{\small KJØBENHAVN}\\
{\scriptsize TRYKT I UNIVERSITETSBOGTRYKKERIET (J. H. SCHULTZ)}\\
{\scriptsize 1902}
\end{center}
\clearpage

% --- front matter: title (PDF 8) ---
\begin{center}
{\Huge FILOSOFISKE}\\[10pt]
{\Huge PROBLEMER}\\[1.5em]
{\scriptsize AF}\\[1.5em]
{\large HARALD HØFFDING}\\[3em]
\rule{1.5cm}{0.4pt}\\[3em]
KJØBENHAVN\\
{\scriptsize TRYKT I UNIVERSITETSBOGTRYKKERIET (J. H. SCHULTZ)}\\
{\scriptsize 1902}
\end{center}
\clearpage

% --- front matter: preface, no heading, in italic (PDF 10) ---
\begin{itshape}
Det første Udkast til denne Afhandling laa til Grund for en Række
Forelæsninger, jeg i Februar Maaned d. A. holdt ved Upsala Universitet efter
Indbydelse fra dettes humanistiske Sektion. Efterat jeg nu har taget de Emner,
jeg behandlede i disse Forelæsninger, op til fornyet Drøftelse, vender min Tanke
sig med hjertelig Hilsen og Tak til den minderige Højskole ved Fyrisaa, som for
en Tid vilde regne mig blandt sine Lærere, --- og til de ældre og yngre Venner,
jeg der vandt.

\bigskip
16. August 1902.

\begin{flushright}HARALD HØFFDING.\end{flushright}
\end{itshape}
\clearpage

% --- front matter: contents (PDF 12) ---
\begin{center}{\large INDHOLD}\end{center}
\begin{flushright}{\scriptsize Side}\end{flushright}
\tocline{}{INDLEDNING}{1}
\tocline{I.}{BEVIDSTHEDSPROBLEMET}{6}
\tocsub{1. Personlighedsbegrebet og den analytiske Psykologi}{6}
\tocsub{2. Diskontinuiteter paa det psykiske Omraade}{12}
\tocsub{3. Psykologi og Fysiologi}{22}
\tocsub{4. Villie og Energi}{25}
\tocline{II.}{ERKENDELSESPROBLEMET}{28}
\tocsub{1. Erkendelsens Arter}{28}
\tocsub{2. Erkendelsens Principer}{31}
\tocsub{3. Kvalitet og Kvantitet}{39}
\tocsub{4. Elementært og idealt Aarsagsbegreb}{43}
\tocsub{5. Subjekt og Objekt}{48}
\tocline{III.}{TILVÆRELSESPROBLEMET}{53}
\tocsub{1. Problem og Metode}{53}
\tocsub{2. Metafysik som Kunst}{57}
\tocsub{3. Første Urfænomen (Enhed eller Mangfoldighed)}{59}
\tocsub{4. Andet Urfænomen (Aand eller Materie)}{62}
\tocsub{5. Tredie Urfænomen (Bestaaen eller Udvikling)}{65}
\tocline{IV.}{VURDERINGSPROBLEMET}{70}
\tocsub{1. Indledning}{70}
\tocplain{\hspace*{1em}\textit{A. Det etiske Problem.}}
\tocplain{\hspace*{3em}a. Det etiske Arbejde.}
\tocsub{2. Kontinuitetsprincipet i Etiken}{72}
\tocplain{\hspace*{3em}b. Den etiske Vurderings Rationalitet.}
\tocsub{3. Grundværdiernes Kamp}{75}
\tocsub{4. De individuelle Betingelsers Forskellighed}{77}
\tocplain{\hspace*{1em}\textit{B. Det religiøse Problem.}}
\tocsub{5. Kontinuitetsprincipet i Religionsfilosofien}{78}
\tocsub{6. Religionens psykologiske Sted}{80}
\tocsub{7. Historiske Former for Religion}{82}

\medskip
\noindent Noter S. 85.
\begin{center}\rule{2cm}{0.4pt}\end{center}
\clearpage

% ===== TEXT (pp. 1--84) =====
%%BODY%%
\clearpage

% ===== NOTER (pp. 85--90) =====
\begin{small}
%%NOTES%%
\end{small}

\end{document}
